\section{TrainerV2}

\subsection{Scopo del TrainerV2}

Il \texttt{TrainerV2} è una versione avanzata del trainer base del progetto.
Non modifica l'architettura del detector, il \textit{Target Generator}, la
funzione di loss o il processo di post-processing, ma modifica principalmente
la strategia utilizzata per ottimizzare i pesi del modello.

La classe eredita direttamente da \texttt{Trainer}:

\begin{verbatim}
class TrainerV2(Trainer):
\end{verbatim}

e riutilizza quindi la maggior parte della logica già presente nel trainer
base, come la gestione dei dataset, la costruzione dei target, il calcolo della
loss, il gradient clipping, la validazione, il logging e i checkpoint.

Le principali modifiche introdotte da \texttt{TrainerV2} riguardano
l'ottimizzazione:

\begin{itemize}
    \item iniziale congelamento della ResNet;
    \item \textit{learning rate warmup};
    \item \textit{cosine annealing};
    \item learning rate differente per backbone e detection layers;
    \item \textit{Exponential Moving Average} (EMA) dei pesi;
    \item gradient clipping.
\end{itemize}

La struttura generale rimane quindi:

\[
\text{input}
\rightarrow
\text{model}
\rightarrow
\text{prediction}
\rightarrow
\text{target}
\rightarrow
\text{loss}
\rightarrow
\text{backpropagation}
\rightarrow
\text{optimizer}.
\]

Le modifiche del TrainerV2 riguardano soprattutto il modo in cui
l'optimizer aggiorna i parametri del modello nel tempo.

\subsection{Freeze e Unfreeze del Backbone}

Il backbone del detector è costituito dalla ResNet-50. Nel progetto è
possibile mantenerlo congelato per un certo numero di epoche iniziali:

\[
\text{ResNet frozen}
\]

Durante questa fase i parametri della ResNet hanno:

\begin{verbatim}
requires_grad = False
\end{verbatim}

e quindi non vengono aggiornati durante il backpropagation.

Al contrario, FPN e detection heads rimangono allenabili.

La logica è:

\[
\text{ResNet}
\rightarrow
\text{frozen}
\]

\[
\text{FPN + Heads}
\rightarrow
\text{trainable}.
\]

L'idea è sfruttare inizialmente le feature già apprese dal backbone
pretrained e concentrarsi sull'apprendimento delle parti specifiche del
detector.

Dopo il numero di epoche specificato, la ResNet viene sbloccata:

\[
\text{ResNet frozen}
\rightarrow
\text{ResNet trainable}.
\]

Nel progetto, quando la ResNet viene congelata, viene anche mantenuta in
modalità evaluation:

\begin{verbatim}
self.backbone.eval()
\end{verbatim}

in modo da evitare che le statistiche interne della Batch Normalization
cambino durante la fase di congelamento.

Il congelamento del backbone viene quindi utilizzato come una forma di
\textit{staged training}: inizialmente si adattano le parti nuove della rete e
successivamente si permette anche al backbone di adattarsi al nuovo dominio.

\subsection{Learning Rate Warmup}

Il \textit{warmup} è una tecnica utilizzata per evitare di iniziare il
training con un learning rate troppo elevato.

Il learning rate determina la dimensione degli aggiornamenti applicati ai
parametri:

\[
\theta_{t+1}
=
\theta_t
-
\eta_t
\nabla_\theta L.
\]

Se il learning rate $\eta_t$ è troppo grande nelle prime iterazioni,
gli aggiornamenti possono essere troppo aggressivi e rendere il training
instabile.

Nel TrainerV2 il learning rate viene quindi inizialmente impostato a un valore
molto basso e aumentato progressivamente fino al learning rate base.

Nel codice viene utilizzato un \texttt{LinearLR}, con:

\[
\text{start factor}=0.01
\]

e:

\[
\text{end factor}=1.
\]

Con un learning rate base $\eta_0$, durante il warmup il learning rate cresce
quindi progressivamente da:

\[
0.01\eta_0
\]

fino a:

\[
\eta_0.
\]

Il processo può essere rappresentato come:

\[
\text{small LR}
\rightarrow
\text{progressive increase}
\rightarrow
\text{base LR}.
\]

Nel progetto il warmup dura due epoche per impostazione predefinita.

L'obiettivo non è migliorare direttamente la capacità del modello, ma rendere
più controllati i primi aggiornamenti dei pesi.

\subsection{Cosine Annealing}

Dopo il warmup, il TrainerV2 utilizza un
\textit{Cosine Annealing Learning Rate Scheduler}.

L'idea è far diminuire gradualmente il learning rate durante il training
secondo una funzione coseno.

Una formulazione semplificata è:

\[
\eta_t
=
\eta_{\min}
+
\frac{1}{2}
(\eta_{\max}-\eta_{\min})
\left(
1+\cos\left(
\frac{t}{T}\pi
\right)
\right).
\]

Dove:

\begin{itemize}
    \item $\eta_{\max}$ è il learning rate iniziale;
    \item $\eta_{\min}$ è il learning rate minimo;
    \item $t$ è il punto corrente del training;
    \item $T$ è la durata del ciclo.
\end{itemize}

All'inizio:

\[
\eta_t \approx \eta_{\max},
\]

mentre verso la fine:

\[
\eta_t \rightarrow \eta_{\min}.
\]

L'idea intuitiva è utilizzare aggiornamenti più grandi quando il modello è
ancora lontano da una buona soluzione e ridurre progressivamente la loro
dimensione quando il training si avvicina alla regione di convergenza.

Nel TrainerV2 il warmup e il cosine annealing sono concatenati:

\[
\text{Warmup}
\rightarrow
\text{Cosine Annealing}.
\]

Il learning rate segue quindi inizialmente una crescita controllata e
successivamente una decrescita graduale.

\subsection{Differential Learning Rate}

Quando la ResNet viene sbloccata, il TrainerV2 non utilizza necessariamente lo
stesso learning rate per tutte le parti del modello.

I parametri vengono infatti divisi in due gruppi:

\[
\text{Backbone}
\]

e

\[
\text{FPN + Detection Heads}.
\]

Nel progetto il learning rate del backbone è impostato a una frazione di quello
utilizzato dalle parti nuove:

\[
\eta_{\text{backbone}}
=
0.1\eta_{\text{base}}.
\]

Quindi, con:

\[
\eta_{\text{base}}=10^{-4},
\]

si ottiene:

\[
\eta_{\text{backbone}}=10^{-5}
\]

mentre FPN e detection heads continuano a utilizzare:

\[
\eta_{\text{head}}=10^{-4}.
\]

La motivazione è che la ResNet è già pretrained e contiene feature
general-purpose, mentre FPN e detection heads sono state introdotte per il
problema specifico.

È quindi ragionevole voler modificare il backbone più lentamente rispetto
alle parti nuove:

\[
\text{pretrained features}
\rightarrow
\text{small updates}
\]

\[
\text{new detector layers}
\rightarrow
\text{larger updates}.
\]

Questa strategia viene spesso indicata come \textit{differential learning
rate}.

\subsection{Exponential Moving Average}

Il TrainerV2 mantiene inoltre una seconda rappresentazione dei parametri del
modello tramite una \textit{Exponential Moving Average} (EMA).

L'idea è non utilizzare esclusivamente i pesi del modello nell'ultimo
optimization step, ma mantenere una media mobile dei pesi osservati durante il
training.

Se $\theta_t$ sono i parametri correnti e $\theta_t^{EMA}$ i parametri mediati,
l'aggiornamento è:

\[
\theta_t^{EMA}
=
\alpha\theta_{t-1}^{EMA}
+
(1-\alpha)\theta_t.
\]

Dove $\alpha$ è il fattore di decadimento.

Nel progetto:

\[
\alpha=0.999.
\]

Un valore elevato significa che la EMA cambia lentamente e quindi tende a
smussare le variazioni rapide dei parametri.

In forma intuitiva:

\[
\text{EMA}
=
\text{versione più stabile dei pesi nel tempo}.
\]

La EMA non viene utilizzata al posto del modello durante il normale
backpropagation. Il modello continua infatti ad essere aggiornato normalmente
dall'optimizer.

La copia EMA viene invece aggiornata dopo ogni:

\[
\text{optimizer.step()}.
\]

Durante la validazione, il TrainerV2 sostituisce temporaneamente i pesi del
modello con quelli EMA:

\[
\text{current weights}
\rightarrow
\text{EMA weights}
\rightarrow
\text{validation}.
\]

Terminata la validazione, i pesi originali vengono ripristinati.

In questo modo si può valutare il modello utilizzando una versione più
smussata dei parametri senza alterare lo stato reale del training.

\subsection{Gradient Clipping}

Il TrainerV2 mantiene inoltre il \textit{gradient clipping} già presente nel
Trainer base.

Dopo il backward:

\[
L
\rightarrow
\nabla_\theta L
\]

viene limitata la norma complessiva dei gradienti.

Nel progetto:

\begin{verbatim}
torch.nn.utils.clip_grad_norm_(
    self.model.parameters(),
    max_norm=1.0
)
\end{verbatim}

La norma dei gradienti viene quindi limitata a:

\[
\|\nabla_\theta L\|\leq1.
\]

Se la norma supera il limite, i gradienti vengono ridimensionati in modo da
rientrare nel valore massimo.

Questo è particolarmente utile quando il training produce gradienti molto
grandi, perché impedisce che un singolo aggiornamento provochi una variazione
eccessiva dei parametri.

Il gradient clipping non modifica direttamente la loss: modifica il
\textit{gradient} utilizzato dall'optimizer per aggiornare i parametri.

\subsection{Training Loop}

Il ciclo di training del TrainerV2 segue una sequenza ben definita per ogni
batch.

\[
\text{Images}
\rightarrow
\text{Forward}
\rightarrow
\text{Targets}
\rightarrow
\text{Loss}
\rightarrow
\text{Backward}
\rightarrow
\text{Gradient Clipping}
\rightarrow
\text{Optimizer Step}
\rightarrow
\text{EMA Update}.
\]

In dettaglio, prima le immagini vengono inviate al modello:

\[
x
\rightarrow
f_\theta(x).
\]

Il modello produce quindi le prediction della detection head.

Successivamente il \textit{Target Generator} costruisce i target a partire
dalle bounding box del dataset.

Prediction e target vengono quindi passati alla funzione di loss:

\[
L=
L_{\text{cls}}
+
L_{\text{reg}}
+
L_{\text{ctr}}.
\]

A questo punto viene eseguito il backward:

\[
L
\rightarrow
\nabla_\theta L.
\]

I gradienti vengono limitati tramite gradient clipping e infine
l'optimizer aggiorna i pesi:

\[
\theta_{t+1}
=
\operatorname{Optimizer}
(\theta_t,\nabla_\theta L).
\]

Se EMA è abilitata, dopo l'aggiornamento viene eseguito anche:

\[
\theta_t^{EMA}
=
\alpha\theta_{t-1}^{EMA}
+
(1-\alpha)\theta_t.
\]

Questo processo viene ripetuto per tutti i batch dell'epoca.

\subsection{Validation}

Al termine di ogni epoca il TrainerV2 esegue la validazione.

Durante questa fase:

\[
\text{model.train()}
\rightarrow
\text{model.eval()}.
\]

Se EMA è abilitata, i pesi EMA vengono applicati temporaneamente.

La pipeline diventa:

\[
\text{Validation images}
\rightarrow
\text{Detector}
\rightarrow
\text{Post-processing}
\rightarrow
\text{Detections}
\rightarrow
\text{Metrics}.
\]

Vengono quindi calcolate le metriche del dataset, tra cui AP e AR.

La metrica utilizzata come riferimento principale per selezionare il modello
migliore è l'AP:

\[
\text{best model}
=
\arg\max_{\text{epoch}} AP.
\]

\subsection{Checkpoint}

Il TrainerV2 salva lo stato dell'addestramento in checkpoint.

Un checkpoint non contiene solamente i pesi del modello, ma può contenere
anche:

\begin{itemize}
    \item stato dell'optimizer;
    \item stato del learning rate scheduler;
    \item epoca corrente;
    \item global step;
    \item migliore AP ottenuta;
    \item stato della EMA.
\end{itemize}

Questo permette di interrompere e successivamente riprendere il training
mantenendo non solo i pesi del modello, ma anche lo stato dell'ottimizzazione.

Il TrainerV2 salva inoltre il modello con la migliore AP osservata e il
checkpoint relativo all'ultimo stato del training.

\subsection{Schema complessivo}

Il comportamento del TrainerV2 può essere riassunto nel seguente schema:

\[
\boxed{
\begin{array}{c}
\text{Initialize model} \\
\downarrow \\
\text{Freeze ResNet} \\
\downarrow \\
\text{Warmup} \\
\downarrow \\
\text{Cosine Annealing} \\
\downarrow \\
\text{Unfreeze ResNet} \\
\downarrow \\
\text{Differential Learning Rate} \\
\downarrow \\
\text{Training} \\
\downarrow \\
\text{Gradient Clipping} \\
\downarrow \\
\text{Optimizer Update} \\
\downarrow \\
\text{EMA Update} \\
\downarrow \\
\text{Validation with EMA} \\
\downarrow \\
\text{Checkpoint}
\end{array}
}
\]

Nel complesso, quindi, il TrainerV2 non introduce un nuovo detector, ma una
strategia di ottimizzazione più sofisticata rispetto al Trainer base. Le
scelte principali sono finalizzate a rendere più controllato e stabile
l'adattamento di un backbone pretrained alle parti del detector specifiche del
problema.